\documentclass[12pt]{article}

\usepackage{amsthm}

\newtheorem{theorem}{Theorem}
\newtheorem{lemma}{Lemma}
\newtheorem{sublemma}{Lemma}[lemma]

\begin{document}

\section{Proof of Correctness of Table Filling Algorithm}

Let $\delta(x)$ give the set of neighbors of $x$ in $G$. $\oplus$ is the logical xor operation.

\subsection{Proof of Construction of Minimal Digraph}

\begin{lemma}
\label{thm:tbl_correctness1}
$\forall w \forall p \forall q . w \in L_G(p) \oplus w \in L_G(q) \rightarrow$ The table filling algorithm finds $p$ and $q$
to be distinguished
\end{lemma}
\begin{proof}
Proof by strong induction on the length of $w$.

Base case $|w|=1$. Then the table filling algorithm
would have distinguished $p$ from $q$ during the initialization of
the table in the table filling algorithm.

We now show the inductive step. As hypotheses, we have the following:

\begin{itemize}
    \item $w=a_1 a_2 ... a_n$
    \item $w \in L_G(p) \oplus w \in L_G(q)$
    \item $\forall w' \forall p' \forall q' . |w'| < n \rightarrow w' \in L_G(p') \oplus w' \in L_G(q') \rightarrow$ The table filling algorithm finds $p'$ and $q'$
    to be distinguished.
\end{itemize}

Our goal is to show that the table filling algorithm finds $p$ and $q$ to be distinguished.

Consider the following cases:

\begin{enumerate}
    \item The string $w'=a_1$ distinguishes $p$ from $q$. Then we can immediately apply the
    inductive hypothesis to conclude that the table filling algorithm finds $p$ and $q$
    to be distinguished.
    \item The string $w$ distinguishes $p$ from $q$, but the string $w'=a_1$ does not. Then exactly
    one of the following must be true:
    \begin{enumerate}
        \item At least one neighbor $p'$ of $p$ must be distinguished from all neighbors $q'$ of $q$ by $w'=a_2 ... a_n$.
        Consider the following cases:
        \begin{enumerate}
            \item $q$ has no neighbors. Then we will mark and therefore
            distinguish $p$ from $q$ during the first iteration of the algorithm.
            \item $q$ has at least one neighbor. Let $q'$ be one of these neighbors.
            By the inductive hypothesis, the table algorithm finds $p'$ to be
            distinguished from $q'$.
    
            This means that during the iteration of the algorithm, $p'$ and $q'$ will be marked. This allows
            the table filling algorithm to distinguish $p$ from $q$ in the next iteration of the loop.
        \end{enumerate}

        \item At least one neighbor $q'$ of $q$ must be distinguished from all neighbors $p'$ of $p$ by $w'=a_2 ... a_n$.
        We can repeat the same argument given in the previous step without loss of generality.
    \end{enumerate}
\end{enumerate}

\end{proof}

\begin{lemma}
\label{thm:tbl_correctness2}
$\forall w \forall p \forall q.$ The table filling algorithm finds $p$ and $q$ to be distinguished $\rightarrow w \in L_G(p) \oplus w \in L_G(q)$
\end{lemma}
\begin{proof}
Proof by strong induction on the length of $w$.

Base case $|w|=1$. Then during the initialization of the table, we discovered that
the label of $p$ is not equal to the label of $q$. This directly shows that the string consisting
of the label of $p$ is in $L_G(p)$ but not in $L_G(q)$.

We now show the inductive step. As hypotheses, we have the following:

\begin{itemize}
    \item $w=a_1 a_2 ... a_n$
    \item The table filling algorithm $p$ and $q$ to be distinguished.
    \item $\forall w' \forall p' \forall q' . |w'| < n \rightarrow$ The table filling algorithm finds $p'$ and $q'$ to be distinguished $\rightarrow w' \in L_G(p') \oplus w' \in L_G(q')$
\end{itemize}

Our goal is to show that $w \in L_G(p) \oplus w \in L_G(q)$

Consider the following cases:
\begin{enumerate}
    \item The string $w'=a_1$ distinguishes $p$ from $q$. Then we can immediately
    apply the inductive hypothesis to conclude that $w' \in L_G(p) \oplus w' \in L_G(q)$.
    This implies that the $w$, which starts with $a_1$, also distinguishes $p$ and $q$.
    This means that $w \in L_G(p) \oplus w \in L_G(q)$ as desired.
    \item The string $w$ distinguishes $p$ from $q$, but the string $w'=a_1$ does not.
    Then the table filling algorithm must have determined that one of the following is true:
    \begin{enumerate}
        \item At least one neighbor $p'$ of $p$ must have been marked for
        all neighbors $q'$ of $q$. Consider the following cases:
        \begin{enumerate}
            \item $q$ has no neighbors. Then the string $a_1 a_2 \in L_G(p)$ but $a_1 a_2 \notin L_G(q)$.
            Since $w$ starts with $a_1 a_2$, then we have $w \in L_G(p)$ and $w \notin L_G(q)$.
            \item $q$ has at least one neighbor. Let $q'$ be one of these neighbors.
            This means that the algorithm found $p'$ to be distinguished from $q'$.
            By the inductive hypothesis, we then have the following:
            $w' \in L_G(p') \oplus w' \in L_G(q')$ where $w'=a_2 ... a_n$.

            Since $p'$ is a neighbor of $p$ and $q'$ is a neighbor of $q$, this
            directly shows that $w \in L_G(p) \oplus w \in L_G(q')$.
        \end{enumerate}
        \item At least one neighbor $q'$ of $q$ must have been marked for all
        neighbors $p'$ of $p$. We can repeat the same argument from the previous
        step without loss of generality.
    \end{enumerate}
\end{enumerate}

\end{proof}

\begin{theorem}
For all strings $w$ and vertices $p$, $q$, the table filling algorithm
finds $p$ and $q$ to be distinguished if and only if $w$ is in exactly
one of $L_G(p)$ and $L_G(q)$
\end{theorem}
\begin{proof}
This immediately follows by Lemma \ref{thm:tbl_correctness1} and Lemma \ref{thm:tbl_correctness2}.
\end{proof}

\clearpage

\subsection{Proof of Uniqueness of Minimum}

We prove this by constructing an isomorphism. Suppose we have a graph
$G$ minimized to $G'$ and some arbitrary minimal graph $G_{min}$ where
$L(G')=L(G_{min})$. We then show that $G' \cong G_{min}$. This is achieved
by demonstrating the existence of a bijection $f : V(G') \rightarrow V(G_{min})$.

First we note that every node in $G'$ or $G_{min}$ must have a different
language. If this was not the case, then two nodes are indiscernable and
could be merged together. This would contradict the assumption that
both these graphs are minimal.

\begin{lemma}
$f$ is defined everywhere
\end{lemma}
\begin{proof}
Consider some vertex $p$ in $G$. Let $S=L_{G'}(p)$. Then $S \in L(G')$. However
$L(G')=L(G_{min})$, so $S \in L(G_{min})$. This means there must be a vertex $p'$
such that $S=L_{G_{min}}(p')$. We can then map $p \mapsto p'$ in $f$.
\end{proof}

\begin{lemma}
$f$ is well defined
\end{lemma}
\begin{proof}
$f$ is well defined if $p=q$ implies $f(p)$ = $f(q)$.

Assume $f$ is not well defined. Then there is some $p$ such that
$f(p)=p'$ and $f(p)=p''$ where $p' \neq p''$.

This means that $L_{G_{min}}(p') = L_{G_{min}}(p'')$ which means that $p'$
and $p''$ are indistinguishable. This contradicts the assumption that
$G_{min}$ is minimal.
\end{proof}

\begin{lemma}
$f$ is injective
\end{lemma}
\begin{proof}
Assume that $f$ is not injective. Then there are some states $p$ and $q$
such that $f(p)=f(q)$.

Since $L(G)=L(G_{min})$, this implies that $p$ and $q$ are indiscernable.
However our algorithm verified that $p$ and $q$ are discernable, which
is a contradiction.
\end{proof}

\begin{lemma}
$f$ is surjective
\end{lemma}
\begin{proof}
Assume $f$ is not surjective. Then there is some state $p'$ such that there
is no node $p$ such that $f(p)=p'$. This implies that $L(G') \neq L(G_{min})$,
which is a contradiction.
\end{proof}

\end{document}

